\chapter{Abstract}

Wayclip is an ecosystem of tools providing Linux users -- and in the future Windows users -- a new way to create, manage and upload clips.
Wayclip was initially developed specifically for Wayland, however, is now growing to more distributions and compositors.

Wayclip solves the core problem of the lack of open-source clipping utilities for gamers on Linux that provide an easy, dedicated and an interactive way to make and share clips.
Wayclip achieves this by providing a custom daemon in addition to a range of tools.
These tools include a CLI, the Wayclip Application and most notably the Wayclip Platform, enabling the sharing of clips online and engagement between communities.
The Wayclip daemon is analogous to other popular applications, such as NVIDIA ShadowPlay, AMD ReLive, and more, but works by interacting directly with your compositor to pipe the output from PipeWire into GStreamer to be processed by your graphics card.

What makes Wayclip different from other projects is our dedication to the open-source community, by providing a lightweight and ready-to-use solution whilst protecting your privacy.
This project has a small group of developers, which allows us to perfect and express our ideas more clearly, making Wayclip a better project.
Everyday, we are working towards making Wayclip a more organised, more efficient and a more user-friendly environment.

This paper will focus on providing a more detailed outline on the features mentioned, as well as our goal and current development status.
